\scanpage{7}
\uline{Cor} $\mathcal A_1\subset S_1(k)$ stable par \unsure{primitivation} canonique (i.e. celle qui s'annule à l'origine), contient $\log(1+t_1)$, $\exp t_1$, $\cos t_1$, $\sin t_1$, $\operatorname{tg}t_1$ $\ldots$

Soit $X$ schéma. formel lisse de dim $n=p+q$. On sait ce que c'est qu'un \uline{feuilletage de $X$}\inserted{de dim $p$}, i.e. une variété formelle quotient $B$ de $X$, avec $X\unsure{-|\to}B$ submersif, $\dim B=q=n-p$. A un tel feuilletage $\varphi$ est associé un sous-module $\Phi\subset\Theta(X)$ [champs de vecteurs le long des fibres] qui est un ``système de Pfaff de dim $p$ complètement intégrable'', et réciproquement. Si $X$ est un $\mathcal A$-schéma, on dit qu'on a un $\mathcal A$-feuilletage si $\exists$\inserted{sur} $B$ une $\mathcal A$-structure (d'ailleurs unique\unsure{\}|)} grâce à f) telle que $X\longrightarrow B$ soit un $\mathcal A$-morphisme. Cela \uline{implique} que le sous-module $\Phi\subset\Theta(X)$ (ou sous-fibré de $T(X)$) soit un $\mathcal A$-sous-fibré -- i.e. qu'on ait un $\mathcal A$-syst. de Pfaff. compl. intégrable. On veut qu'on ait une réciproque, savoir

h) Si $\Phi_{\mathcal A}$ est un $\mathcal A$-syst. de Pfaff compl. intégrable sur un $\mathcal A$-schéma. formel $X$, alors il donne lieu à un $\mathcal A$-feuilletage \deleted{\unsure{compl. intégrable}}.
\margin{\uline{NB} h) implique g), qui équivaut à h) pour $p=1$}
\figure{\begin{tikzpicture}[x=1cm,y=1cm]\draw (0,2.4)..controls(1.1,3.2)and(1.6,2.4)..(.9,1.7)..controls(.5,1.2)and(.5,.5)..(0,.2);\end{tikzpicture}}
[En d'autres termes, si $i:Y\hookrightarrow X$ est un $\mathcal A$-morphisme donné tel que le composé
\[
T(Y)_0\xrightarrow{T(i)_0}T(X)_0\unsure{-|\to}T(X)_0/\Phi(0)
\]
soit un isom, alors l'unique morphisme
\[
\varphi:X\unsure{-|\to}Y
\]
tel que $\varphi\circ i=\operatorname{id}_Y$ et que $\varphi$ soit compatible avec $\Phi$, est un $\mathcal A$-morphisme $\ldots$

h') la fibre de $0$ dans un tel feuilletage est un sous-$\mathcal A$-variété formelle$\ldots$]
\figure{\begin{tikzpicture}[x=1cm,y=1cm]\draw (0,.8)..controls(-.2,.1)and(-.3,-.1)..(0,-.15);\draw (.05,.35)..controls(.2,.2)and(.6,.3)..(.6,.15)..controls(.5,0)and(.2,.1)..(.05,.15);\end{tikzpicture}}
\margin{6}
